% Part: methods
% Chapter: induction
% Section: introduction

\documentclass[../../../include/open-logic-section]{subfiles}

\begin{document}

\olfileid[pt]{mth}{ind}{int}

\olsection{Introdução}

A indução é uma técnica de prova importante, usada, em diferentes
formas, em quase todas as áreas da lógica, da ciência da computação teórica e da
matemática.  Ela é necessária para provar muitos dos resultados em lógica.

A indução é frequentemente contrastada com a dedução, e caracterizada como a
inferência do particular para o geral.  Por exemplo, se
observamos muitas esmeraldas verdes, e nada que chamaríamos de esmeralda
que não seja verde, poderíamos concluir que todas as esmeraldas são verdes. Essa
é uma inferência indutiva, no sentido de que procede de muitos casos
particulares (esta esmeralda é verde, aquela esmeralda é verde, etc.) a uma
afirmação geral (todas as esmeraldas são verdes).  A indução \emph{matemática}
também é uma inferência que conclui uma afirmação geral, mas é de um
tipo muito diferente dessa ``indução simples.''

De modo bem aproximado, uma prova indutiva em matemática conclui que todos os
objetos matemáticos de certo tipo têm uma certa propriedade.  No
caso mais simples, os objetos matemáticos com que uma prova indutiva se
ocupa são os números naturais.  Nesse caso, uma prova indutiva
é usada para estabelecer que todos os números naturais têm alguma propriedade, e
faz isso mostrando que
\begin{enumerate}
    \item $0$ tem a propriedade, e
    \item sempre que um número~$k$ tem a propriedade, também~$k+1$ a tem.
\end{enumerate}
A indução sobre os números naturais pode então, também, ser frequentemente usada para provar
afirmações gerais sobre objetos matemáticos aos quais se podem atribuir números. Por
exemplo, conjuntos finitos têm, cada um, um número finito~$n$ de elementos, e
se pudermos usar a indução para mostrar que todo número~$n$ tem a propriedade
``todos os conjuntos finitos de tamanho~$n$ são \dots'' então teremos mostrado
algo sobre todos os conjuntos finitos.

A indução também pode ser generalizada para objetos matemáticos que são
\emph{definidos indutivamente}.  Por exemplo, as expressões de uma linguagem
formal, como as da lógica de primeira ordem, são definidas indutivamente.
A \emph{indução estrutural} é um modo de provar resultados sobre todas essas
expressões.  A indução estrutural, em particular, é muito
útil---e amplamente usada---em lógica.

\end{document}
